\clearpage
\fancyhead[L]{\small JUNO.}
\hypertarget{native-268-lower}{}
\begin{center}
JUNO.
\end{center}

\noindent\rule{\linewidth}{0.4pt}

\begin{center}
Monatliche Correspondenz. Band XI. S. 475–482. 1805 Mai.
\end{center}

\noindent\rule{\linewidth}{0.4pt}

Wir liefern hier eine kleine Nachlese von Beobachtungen der \emph{Juno}, die auf entferntern Sternwarten gemacht wurden und erst später zu unserer Kenntniss gelangten. Sie umfassen den letzten Zeitraum ihrer diesmaligen Sichtbarkeit, und konnten zum Theil schon nicht mehr im Meridian, sondern nur an Aequatorial-Instrumenten ausserhalb desselben gemacht werden. Vorzüglich war der helle Mondschein, der in den ersten Tagen des Februars Statt fand, Ursache, dass \emph{Juno} nicht mehr gesehen und dann nicht wieder aufgefunden werden konnte. Nur Dr. \textsc{Gauss} war noch am 20. Febr. so glücklich, eine Beobachtung der Juno zu erhalten, die letzte, die uns überhaupt mitgetheilt worden ist.

Aus Palermo von \textsc{Piazzi} erhielten wir noch eine Reihe, an seinem ganzen Kreise gemachter Meridian-Beobachtungen, die wir hier folgen lassen:
\SourceLeaf{269}
\begin{center}
\begingroup\fontsize{9}{11}\selectfont\setlength{\tabcolsep}{3pt}\renewcommand{\arraystretch}{1.1}
\begin{tabular}{r|r|r|r}
1805 & \shortstack{Mittlere Zeit\\in Palermo} & \shortstack{Scheinb. gerade\\Aufsteig. der $\JunoSymbol$} & \shortstack{Scheinb. südl.\\Abweich. der $\JunoSymbol$}\\ \hline
Nov. 22 & $7^{\mathrm u}\ 42^{\mathrm m}\ 39^{\mathrm s}9$ & $357^\circ\ 12'\ 28''5$ & $.\ .\ .\ .\ .$ \\
24 & $7\ 36\ 12.5$ & $357\ 33\ 36.6$ & $10^\circ\ 49'\ 34''9$ \\
25 & $7\ 33\ 1.5$ & $357\ 44\ 54.9$ & $10\ 47\ 21.5$ \\
26 & $7\ 29\ 52.3$ & $357\ 56\ 36.9$ & $10\ 44\ 49.9$ \\
28 & $7\ 23\ 37.0$ & $358\ 20\ 49.5$ & $10\ 38\ 59.5$ \\
29 & $7\ 20\ 31.4$ & $358\ 33\ 26.1$ & $10\ 35\ 43.2$ \\
Dec. 4 & $7\ 5\ 28.2$ & $359\ 42\ 42.3$ & $10\ 15\ 52.4$ \\
10 & $6\ 47\ 21.8$ & $1\ 5\ 11.9$ & $.\ .\ .\ .\ .$ \\
\end{tabular}\endgroup
\end{center}

\par\indent In der Beobachtung vom 10. Decemb. scheint ein Fehler zu liegen, indem sich diese beträchtlich von dem aus den Elementen berechneten Orte entfernt. Sorgfältig verglich Dr. \textsc{Gauss} diese Beobachtungen mit seinen IV. Elementen, und erhielt folgende Resultate:

\begin{center}
\begingroup\fontsize{9}{11}\selectfont\setlength{\tabcolsep}{3pt}\renewcommand{\arraystretch}{1.1}
\begin{tabular}{r|r|r|r|r}
1805 & \shortstack{Berechn. gerade\\Aufsteigung\\der $\JunoSymbol$} & \shortstack{Berechn. südl.\\Abweichung\\der $\JunoSymbol$} & \multicolumn{2}{c}{\shortstack{Fehler der IV. Ele-\\mente des Dr. \textsc{Gauss}\\in AR\qquad in Decl.}}\\ \hline
Nov. 22 & $357^\circ\ 12'\ 11''6$ & $.\ .\ .\ .\ .$ & $-16''9$ & $.\ .\ .$ \\
24 & $357\ 33\ 28.3$ & $10^\circ\ 49'\ 47''3$ & $-8.3$ & $+12''4$ \\
25 & $357\ 44\ 42.4$ & $10\ 47\ 25.9$ & $-12.5$ & $+4.4$ \\
26 & $357\ 56\ 18.3$ & $10\ 44\ 55.0$ & $-18.6$ & $+5.1$ \\
28 & $358\ 20\ 39.6$ & $10\ 39\ 5.0$ & $-9.9$ & $+5.5$ \\
29 & $358\ 33\ 22.3$ & $10\ 35\ 55.8$ & $-3,8$ & $+12.6$ \\
Dec. 4 & $359\ 42\ 27.6$ & $10\ 16\ 9.3$ & $-14.7$ & $+16.9$ \\
10 &  &  &  &  \\
\end{tabular}\endgroup
\end{center}

\par\indent Ohnerachtet diese Abweichungen so gering sind, dass sie kaum eine Verbesserung der IV. Elemente der Juno zu erfordern scheinen, da alle unsere ältern Planeten-Tafeln in gerader Aufsteigung gleiche, vielleicht oft noch beträchtlichere Abweichungen zeigen, so gründete doch der fleissige Dr. \textsc{Gauss} auf diese Beobachtungen und eine, von ihm selbst am 20. Februar erhaltene, folgende verbesserte V. Elemente der Juno:

\Needspace{18\baselineskip}
\begin{center}
\begingroup\fontsize{8}{10}\selectfont\setlength{\tabcolsep}{3pt}\renewcommand{\arraystretch}{1.1}
\begin{tabular}{lr}
Epoche 1805, Meridian von Seeberg . . . . & $42^\circ\ 32'\ 36''$ \\
\multicolumn{2}{l}{$\left.\begin{array}{lr}\text{\AscendingNode\ 1805 . . . . . . . . . . . . . .}&171\ 4\ 15.6\\\text{Sonnenferne . . . . . . . . . . . . .}&233\ 11\ 39\end{array}\right\}\begin{array}{l}\text{siderisch ruhend}\\\text{vorausgesetzt}\end{array}$} \\
Neigung der Bahn . . . . . . . . . . & $13\ 3\ 38$ \\
tägliche tropische Bewegung . . . . . . . & $815''9595$ \\
Jährliche . . . . . . . . . . . . . . & $82^\circ\ 43'\ 45''2$ \\
Excentricität . . . . . . . . . . . . & $0.254236$ \\
Logarithmus der halben Axe . . . . . . & $0.4256078$ \\
\end{tabular}\endgroup
\end{center}

\par\indent Die hier erwähnte Beobachtung von Dr. \textsc{Gauss} war folgende:

\begin{center}
\begingroup\fontsize{9}{11}\selectfont\setlength{\tabcolsep}{3pt}\renewcommand{\arraystretch}{1.1}
\begin{tabular}{c|c|c}
Mittlere Zeit & AR $\JunoSymbol$ & Decl. bor. $\JunoSymbol$\\ \hline
1805, Febr. 20. $7^{\mathrm u}\ 11^{\mathrm m}\ 12^{\mathrm s}$ & $30^\circ\ 27'\ 2''$ & $1^\circ\ 47'\ 48''$ \\
\end{tabular}\endgroup
\end{center}

\par\indent Der aus den IV. Elementen für diese Zeit berechnete Ort der \emph{Juno} war:

\begin{center}
\begingroup\fontsize{9}{11}\selectfont\setlength{\tabcolsep}{3pt}\renewcommand{\arraystretch}{1.1}
\begin{tabular}{c|c|c|c}
& & \multicolumn{2}{c}{Unterschied}\\ AR $\JunoSymbol$ & Decl. bor. $\JunoSymbol$ & in AR & in Decl.\\ \hline
$30^\circ\ 27'\ 41''6$ & $1^\circ\ 47'\ 20''$ & $+39''6$ & $-27''7$ \\
\end{tabular}\endgroup
\end{center}
\SourceLeaf{270}
\par\indent Nach jenen fünften Elementen berechnete Dr. \textsc{Gauss} nachfolgende Ephemeride für den geocentrischen Lauf der Juno [von 1805 October 20 bis 1806 Juli 23 wie für Ceres und Pallas d. B. S. 262. 263].

[In dem übrigen Theil dieses Aufsatzes gibt von \textsc{Zach} noch Beobachtungen von \textsc{Oriani} in Mailand und von \textsc{Olbers} in Bremen so wie die Vergleichung dieser Beobachtungen mit \textsc{Gauss} IV. Elementen und auch die von Dr. \textsc{Burckhardt} berechneten Elemente der Bahn der Juno nach dem \emph{Moniteur} vom 28. Dec. 1804.]

[Die auf Berliner Zeit bezogenen von \textsc{Gauss} berechneten Elemente: die X\textsuperscript{ten} der Ceres die VIII\textsuperscript{ten} der Pallas und die V\textsuperscript{ten} der Juno hat \textsc{Bode} in dem Astronomischen Jahrbuch für 1808 Berlin 1805 S. 270 veröffentlicht.]

\Needspace{12\baselineskip}
\fancyhead[L]{\small ERSTER COMET VON 1805.}
\begin{center}
ERSTER COMET VOM JAHR 1805.

\rule{0.6\linewidth}{0.4pt}

Monatliche Correspondenz. Band XIII. S. 79–83. 1806 Januar.

\rule{0.6\linewidth}{0.4pt}

—\quad—\quad—
\end{center}

Dr. \textsc{Gauss} berichtet uns vom 5. December aus Braunschweig, dass er, wiewohl vergeblich versucht habe, diese \textsc{Bessel}'schen Elemente den Frankfurter und Marseiller Beobachtungen anzupassen; allein diese Beobachtungen scheinen gar wenig genau zu sein. „Da diese Elemente“ schreibt Dr. \textsc{Gauss}, „mehrere Grade von \textsc{Huth}'s letzten Beobachtungen abweichen, so machte ich einen Versuch, sie nach diesen zu verbessern, obgleich \textsc{Huth}'s Angaben, die nur in runden Fünfern von Minuten gesetzt sind, blos Ocular-Schätzungen zu sein scheinen; dieses wollte aber nicht gelingen. Auch die mir von Ihnen gütigst mitgetheilten \textsc{Thulis}'schen Beobachtungen (November-Heft 1805 S. 502) scheinen nicht sonderlich zu sein, wenigstens habe ich auch damit zu keinem befriedigenden Resultate gelangen können; man sieht dies auch schon aus der Vergleichung mit \textsc{Olbers}' Beobachtungen, die am 29. October um sechs Minuten in gerader Aufsteigung und neun Min. in der Abweichung, am 31. October um eine Min. in gerader Aufsteigung und zwölf Minuten in der Abweichung mit diesen differiren, wenn man sie auf einerlei Zeit reducirt; indessen weichen \textsc{Bessel}'s Elemente am 9. November um 71 Minuten in der Länge von \textsc{Thulis}' Beobachtungen ab, welches doch mehr ist, als man Fehler der Beobachtung voraussetzen darf. Ich habe daher einen Versuch gemacht, den Fehler zu corrigiren, habe aber doch einen Fehler von beinahe 29 Minuten in der Breite nicht vermeiden können. Diese sind die Elemente

\begin{center}
\begingroup\fontsize{8}{10}\selectfont\setlength{\tabcolsep}{3pt}\renewcommand{\arraystretch}{1.1}
\begin{tabular}{lr}
Durchgang durch das Perihelium November . . . . & $17.746$ in Seeberg. \\
Länge des aufsteigenden Knotens . . . . . . . . & $340^\circ\ 11'$ \\
Neigung der Bahn . . . . . . . . . . . . . & $17\ 34$ \\
Länge der Sonnennähe . . . . . . . . . . . . & $157\ 17$ \\
Logarithmus des kleinsten Abstandes . . . . . . . & $9.53969$ \\
Bewegung . . . . . . . . . . . . . . . . & rechtläufig \\
\end{tabular}\endgroup
\end{center}

„Diese Resultate verdienen indessen gleichfalls wenig Vertrauen.“
\SourceLeaf{271}
\par\indent Wir haben alle diese Umstände blos deshalb angeführt, um andern Berechnern eine vergebliche Mühe zu ersparen, im Fall sie es versuchen wollten, die Frankfurter und Marseiller Beobachtungen zu irgend einer Übereinstimmung zu bringen. Da die \textsc{Bessel}'sche Bahn die \textsc{Olbers}'schen Beobachtungen bis zum 13. November so befriedigend darstellt, so dürfen diese Elemente als genau entwickelt angesehen werden, und wenn wir künftig in Besitz besserer und längerer Beobachtungen kommen sollten, nur sehr geringe Verbesserungen erleiden.

\hfill[\textsc{von Zach.}]

\Needspace{14\baselineskip}
\fancyhead[L]{\small ZWEITER COMET VON 1805.}
\begin{center}
ZWEITER COMET VOM JAHR 1805.

\rule{0.6\linewidth}{0.4pt}

Monatliche Correspondenz. Band XIII. S. 83–91. 1806 Januar.

\rule{0.6\linewidth}{0.4pt}

—\quad—\quad—
\end{center}

Dr. \textsc{Gauss} beobachtete diesen Cometen den 8. December; er wurde mit 359 im Wassermann und einem andern Stern der \emph{Hist. célest.} verglichen.

\begin{center}
\begingroup\fontsize{9}{11}\selectfont\setlength{\tabcolsep}{3pt}\renewcommand{\arraystretch}{1.1}
\begin{tabular}{r|r|r|r}
1805 & \shortstack{Mittl. Zeit in\\Braunschweig} & \shortstack{Scheinb. ger.\\Aufsteigung} & \shortstack{Scheinb. südl.\\Abweichung}\\ \hline
December 8 & $6^{\mathrm u}\ 55^{\mathrm m}\ 17^{\mathrm s}$ & $353^\circ\ 7'\ 40''$ & $23^\circ\ 36'\ 24''$ \\
8 & $7\ 58\ 3$ & $352\ 57\ 57$ & $24\ 6\ .\ .$ \\
\end{tabular}\endgroup
\end{center}

\par\indent Die zweite Declination ist blos Schätzung. Aus dieser Beobachtung, der ersten \textsc{Bouvard}'schen vom 16. November, und der obigen des Dr. \textsc{Olbers} vom 2. December berechnete Dr. \textsc{Gauss} folgende parabolische Elemente:

\begin{center}
\begingroup\fontsize{8}{10}\selectfont\setlength{\tabcolsep}{3pt}\renewcommand{\arraystretch}{1.1}
\begin{tabular}{lr}
Durchgang durch das Perihelium den 31. December 1805 , & $7^{\mathrm u}\ 20^{\mathrm m}\ 39^{\mathrm s}$ Seeberger Zeit \\
Länge des Periheliums , . . . . . . . . . . . . & $109^\circ\ 23'\ 40''$ \\
Länge des aufsteigenden Knotens . . . . . . . . . . & $250\ 33\ 14$ \\
Neigung der Bahn . , . . . . . \raisebox{.3ex}{.} . . . . . . . & $16\ 33\ 33$ \\
Logarithmus des kleinsten Abstandes . . . . . . . . & $9.9502477$ \\
Bewegung . . . . . . . . . . . . . . . . . . & rechtläufig. \\
\end{tabular}\endgroup
\end{center}

\par\indent Nach der Methode des Aufsatzes im Mai-Heft der M. C. 1804 [Seite 94 und 120 d. B.] werden die Coordinaten durch folgende Formeln dargestellt:

\[
\begin{gathered}
x = \frac{\alpha \sin(v + 198^\circ\,37'\,20'')}{\cos\frac{1}{2}v^2},\\[4pt]
y = \frac{\ObservedBeta \sin(v + 103^\circ\,41'\,12'')}{\cos\frac{1}{2}v^2},\\[4pt]
z = \frac{\gamma \sin(v + 148^\circ\,41'\,54'')}{\cos\frac{1}{2}v^2}
\end{gathered}
\]

\noindent wo $v$ die Anomalie des Cometen bedeutet, und Log. $\alpha = 9.933968$, Log. $\ObservedBeta = 9.930525$, Log. $\gamma = 9.551021$.

Den Abstand des Cometen von der Erde findet Dr. \textsc{Gauss} nach seinen Elementen mit Rücksicht auf die Parallaxe Nov. 16: 0.24990, Dec. 2: 0.08913, Dec. 8 I. Beob.: 0.04850, Dec. 8 II. Beob.: 0.04841.

Hingegen die Lichtstärke, wenn man die in der Distanz 1. von der Erde und Sonne zur Einheit annimmt November 16: 16.6, December 2: 117.8, December 8: 439.2.

\begin{center}
—\quad—\quad—
\end{center}

\hfill[\textsc{von Zach.}]
